% =============================================================================
% s8_discussion.tex --- §8 Discussion and Limitations (S version, ICLR 9-page).
% Compressed from drafts/s8_discussion.tex (1580 words) to <=300 words / ~0.5 page.
% Reflection-style limitations: "evaluation bounded by X; Y is an open question
% listed as future work" --- honest, not apologetic, not buried (STORY drift #10).
% Numbers carried verbatim from the long draft (locked); none invented (R6).
% Red lines: R3/R11 (no global triage win), R8 (diffusion only as rejected
% contrast), R12 (GradProm weld), R1/R2 (no absolutes / "we derive").
% =============================================================================

\section{Discussion and Limitations}\label{sec:discussion}

\paragraph{Relation to gradient-repair (GradProm).}
GradProm~\citep{gradprom2025} also observes that enhancement can harm diagnosis in
the ISIC domain, but treats the symptom by repairing gradients during training. Our
contribution is orthogonal: we (i)~\emph{derive} a diagnosis-preserving information
lower bound (\cref{lem:dp}), and (ii)~add a \emph{routing} decision
with a query-for-retake channel. Where GradProm fixes gradients, we bound the
information loss \emph{and} decide when not to enhance at all---complementary,
not competing.

\paragraph{Why generative enhancement is rejected.}
We avoid diffusion- and GAN-based restoration: optimising perceptual realism in
dermoscopy can synthesise plausible-but-absent structure (pigment networks, vessels),
a diagnostically misleading medical-safety hazard. Deterministic restoration
trades some sharpness for an information-preservation property a generative model
cannot guarantee; diffusion enters only as a rejected cautionary contrast.

\paragraph{What the method establishes, and its honest boundary.}
The positive contribution is concrete and statistically attributable: within the
moderate-degradation window the diagnosis-preserving objective lets \visienhance retain the
discriminative information a frozen probe relies on (DP-Loss ablation
$\Delta\mathrm{AUC}_{\text{enh}}=+0.0205$, $95\%$ CI $[+0.005,+0.035]$; McNemar
$p=2.3\times10^{-45}$), and it preserves diagnosis better than six restoration baselines on every one.
We are equally explicit about where the method stops. We make no claim that quality-aware triage beats
direct diagnosis \emph{globally}: at a fixed threshold, direct diagnosis attains the best auto-handled
sensitivity ($0.818$ versus $0.788$ for the strongest variant at a $20\%$ referral budget) with
net-benefit CIs overlapping (\cref{sec:exp-triage}). The positive claim is
local---the enhancer preserves discriminative information within the
moderate-degradation window---and the decision surface (\cref{sec:exp-surface}) shows
recoverability is dimension- and severity-dependent, exactly what the query-for-retake
gate is built to exploit.

\paragraph{Limitations and future work.}
Several limitations are themselves the motivation for the retake gate, stated plainly.
(1)~The theory is derived under explicit assumptions and is argued rather than
watertight; \cref{thm:agent-compact} is band-local and does not aggregate into a
global ordering.
(2)~Degradations are semi-synthetic; prospective validation on genuinely degraded
clinical captures remains open.
(3)~Routing depends on a per-input quality scalar whose annotation and re-fitting cost
we do not eliminate (claims restricted to dermatology, no cross-modality universality).
(4)~We report no clinical deployment; all clinical statements rest on decision-curve and
triage simulation, and diagnosis-preservation results are tied to a single frozen probe.
(5)~Enhancement is uneven across axes: it backfires on contrast and underperforms on
colour-shift (\cref{sec:exp-ablations})---honest failure modes that trigger the
caution/retake channels.
(6)~The malignant-salvage gap is unresolved: enhancement salvages only $5.2\%$ of
misclassified melanomas (a net loss of malignant cases), unchanged by a mask-weighted
retrain; closing it---via diagnosis-aware reconstruction or learned routing
thresholds---is the central open question. These negative results are genuine evidence
that the retake gate is a necessary safety mechanism, not a project shortcoming.
